\hypertarget{faster-cpython-the-specializing-interpreter}{%
\chapter{FASTER CPYTHON: THE SPECIALIZING
INTERPRETER}\label{faster-cpython-the-specializing-interpreter}}

\hypertarget{pep-659-specialized-adaptive-interpreter-architecture}{%
\subsection{17.1 PEP 659 Specialized Adaptive Interpreter
Architecture}\label{pep-659-specialized-adaptive-interpreter-architecture}}

Before Python 3.11, the bytecode interpreter loop executed instructions
generically. For instance, a \passthrough{\lstinline!LOAD\_ATTR!}
instruction looked up attributes using hash table lookups every time.
PEP 659 introduced a \textbf{specializing adaptive interpreter}. Rather
than incurring the cost of compiling bytecode to native machine code
(JIT), the interpreter dynamically replaces general opcodes with
optimized, specialized opcodes at runtime based on the actual type of
objects processed.

\hypertarget{instruction-lifecycle-state-machine}{%
\subsubsection{1. Instruction Lifecycle State
Machine}\label{instruction-lifecycle-state-machine}}

Each instruction transitions through a state machine:

\begin{lstlisting}
               [Generic Opcode] (e.g., LOAD_ATTR)
                      |
                      | (First execution)
                      v
               [Adaptive Opcode] (e.g., LOAD_ATTR_ADAPTIVE)
                      |
                      | (Executes N times with same type)
                      v
             [Specialized Opcode] (e.g., LOAD_ATTR_INSTANCE_VALUE)
             /                   \
            / (Type mismatch)     \ (Attribute modification)
           v                       v
     [Deoptimize] ---------> [Generic Opcode]
\end{lstlisting}

\begin{enumerate}
\def\labelenumi{\arabic{enumi}.}
\tightlist
\item
  \textbf{Generic}: Standard instruction compiled from AST.
\item
  \textbf{Adaptive}: When executed, it monitors the types of objects on
  the stack. An internal counter is decremented.
\item
  \textbf{Specialized}: If the counter reaches zero and the operand
  types have been consistent, the opcode is hot-swapped in memory with a
  specialized instruction.
\item
  \textbf{Deoptimized}: If the operand types change (cache miss), the
  opcode executes its slow-path fallback and can revert to the generic
  state.
\end{enumerate}

\begin{center}\rule{0.5\linewidth}{0.5pt}\end{center}

\hypertarget{specialized-bytecode-opcodes-and-inline-cache-memory-layout}{%
\subsection{17.2 Specialized Bytecode Opcodes and Inline Cache Memory
Layout}\label{specialized-bytecode-opcodes-and-inline-cache-memory-layout}}

CPython allocates extra space directly adjacent to the instruction
stream to serve as an \textbf{inline cache}.

\hypertarget{cache-struct-layout-in-c}{%
\subsubsection{1. Cache Struct Layout in
C}\label{cache-struct-layout-in-c}}

In CPython's internal interpreter definitions
(\passthrough{\lstinline!pycore\_code.h!}), the cache entries for
specialized attributes are defined as arrays of structs:

\begin{lstlisting}[language=C]
typedef struct {
    uint16_t counter;           /* Adaptive counter to trigger specialization */
    uint32_t type_version;      /* Version ID of the target object's PyTypeObject */
    uint16_t index;             /* Index of attribute inside object's storage array */
} _PyAttrCache;
\end{lstlisting}

\hypertarget{key-specialized-instructions}{%
\subsubsection{2. Key Specialized
Instructions}\label{key-specialized-instructions}}

\begin{itemize}
\tightlist
\item
  \textbf{\passthrough{\lstinline!LOAD\_ATTR!} -\textgreater{}
  \passthrough{\lstinline!LOAD\_ATTR\_INSTANCE\_VALUE!}}:

  \begin{itemize}
  \tightlist
  \item
    Used when attributes are stored in an instance's values array (split
    dictionary layout).
  \item
    Bypasses dictionary key lookup. It compares the target object's type
    version directly against the cache. If they match, it loads the
    value directly from the cached offset
    \passthrough{\lstinline!index!}.
  \end{itemize}
\item
  \textbf{\passthrough{\lstinline!LOAD\_ATTR!} -\textgreater{}
  \passthrough{\lstinline!LOAD\_ATTR\_SLOT!}}:

  \begin{itemize}
  \tightlist
  \item
    Used when the class defines \passthrough{\lstinline!\_\_slots\_\_!}.
  \item
    Directly retrieves the value from the static struct offset of the
    instance slot, bypassing dictionary lookup entirely.
  \end{itemize}
\item
  \textbf{\passthrough{\lstinline!LOAD\_GLOBAL!} -\textgreater{}
  \passthrough{\lstinline!LOAD\_GLOBAL\_MODULE!} /
  \passthrough{\lstinline!LOAD\_GLOBAL\_BUILTIN!}}:

  \begin{itemize}
  \tightlist
  \item
    Caches the keys version of the module globals dictionary and the
    builtins dictionary.
  \item
    Loads values directly from the stored dict value arrays, avoiding
    dictionary key hashing checks.
  \end{itemize}
\end{itemize}

\begin{center}\rule{0.5\linewidth}{0.5pt}\end{center}

\hypertarget{micro-benchmarking-and-deoptimization-performance-cost}{%
\subsection{17.3 Micro-benchmarking and Deoptimization Performance
Cost}\label{micro-benchmarking-and-deoptimization-performance-cost}}

If code is polymorphic (objects of different classes with varying
attribute offsets flow through the same bytecode instruction), the
specialized instruction will continuously fail validation, triggering a
fallback to the slow-path generic execution. This process is called
\textbf{deoptimization}. Frequent deoptimization ruins performance
because the VM incurs the combined overhead of cache miss checking,
slow-path lookup, and state-machine transitions.

\begin{center}\rule{0.5\linewidth}{0.5pt}\end{center}
